\chapter{Deklination der Nomen}
\label{cha:Deklination}

Die Deklination im Deutschen markiert die grammatische Funktion eines Nomens durch Kasus (Nominativ, Akkusativ, Dativ, Genitiv). Im Gegensatz zu analytischen Sprachen wie dem Englischen behält das Deutsche eine starke flektierende Struktur.

\section{Grundlagen der Kasus}
\label{sec:Kasus}

Die vier Kasus im Deutschen:

\subsection{Nominativ (Wer? Was?)}
\begin{tcolorbox}[title=Funktion]
Der Nominativ bezeichnet das Subjekt des Satzes - wer oder was handelt oder existiert.
\end{tcolorbox}

\begin{itemize}
    \item \textbf{Der} Mann kommt.
    \item \textbf{Das} Kind schläft.
    \item \textbf{Die} Katze ist klein.
    \item \textbf{Ich} bin müde.
    \item \textbf{Sie} arbeitet.
    \item \textbf{Wir} sind fertig.
    \item \textbf{Die} Studenten lernen.
\end{itemize}

\subsection{Akkusativ (Wen? Was?)}
\begin{tcolorbox}[title=Funktion]
Der Akkusativ bezeichnet das direkte Objekt - wen oder was wird direkt betroffen.
\end{tcolorbox}

\begin{itemize}
    \item Ich sehe \textbf{den} Mann.
    \item Sie hat \textbf{das} Buch gelesen.
    \item Er kennt \textbf{die} Frau.
    \item Wir kaufen \textbf{die} Äpfel.
    \item Ich meine \textbf{dich}.
    \item Er liebt \textbf{seine} Frau.
    \item Wir brauchen \textbf{einen} Wagen.
\end{itemize}

\subsection{Dativ (Wem?)}
\begin{tcolorbox}[title=Funktion]
Der Dativ bezeichnet das indirekte Objekt - wem wird etwas gegeben oder zugeordnet.
\end{tcolorbox}

\begin{itemize}
    \item Ich helfe \textbf{dem} Mann.
    \item Sie gibt \textbf{der} Katze Futter.
    \item Er schreibt \textbf{dem} Freund.
    \item Wir danken \textbf{ihr}.
    \item Das gehört \textbf{mir}.
    \item Sie vertraut \textbf{ihm}.
    \item Ich antworte \textbf{dem} Lehrer.
\end{itemize}

\subsection{Genitiv (Wessen?)}
\begin{tcolorbox}[title=Funktion]
Der Genitiv bezeichnet Besitz oder Zugehörigkeit - wessen etwas ist.
\end{tcolorbox}

\begin{itemize}
    \item Das ist das Auto \textbf{des} Mannes.
    \item Die Farbe \textbf{der} Blume ist rot.
    \item Das Haus \textbf{des} Nachbarn ist groß.
    \item Die Mutter \textbf{des Kindes}.
    \item Das Buch \textbf{meines Bruders}.
    \item Die Tür \textbf{des Hauses}.
    \item Der Wagen \textbf{meines Vaters}.
\end{itemize}

\section{Deklinationstabellen}
\label{sec:Deklinationstabellen}

\subsection{Schwache Deklination (bestimmter Artikel)}
\begin{center}
\begin{tabular}{|l|c|c|c|c|}
\hline
\textbf{Kasus} & \textbf{Maskulin} & \textbf{Feminin} & \textbf{Neutral} & \textbf{Plural} \\ \hline
Nominativ & der & die & das & die \\ \hline
Akkusativ & den & die & das & die \\ \hline
Dativ & dem & der & dem & den (+n) \\ \hline
Genitiv & des & der & des & der \\ \hline
\end{tabular}
\end{center}

\subsection{Schwache Deklination (unbestimmter Artikel)}
\begin{center}
\begin{tabular}{|l|c|c|c|c|}
\hline
\textbf{Kasus} & \textbf{Maskulin} & \textbf{Feminin} & \textbf{Neutral} & \textbf{Plural} \\ \hline
Nominativ & ein & eine & ein & (keine) \\ \hline
Akkusativ & einen & eine & ein & (keine) \\ \hline
Dativ & einem & einer & einem & (keinen) \\ \hline
Genitiv & eines & einer & eines & (keiner) \\ \hline
\end{tabular}
\end{center}

\section{Übungen zur Deklination}
\label{sec:DeklinationUebungen}

\begin{tcolorbox}[title={Aufgabe}]
Bestimmen Sie den Kasus und bilden Sie den korrekten Artikel.
\end{tcolorbox}

\begin{enumerate}[label=\textbf{Aufgabe \arabic*}.]
    \item Ich warte auf \underline{\hspace{1.5cm}} Bus. \quad \textit{(Lösung: den -- Akkusativ)}
    \item Das Auto \underline{\hspace{1.5cm}} Nachbarn ist neu. \quad \textit{(Lösung: des -- Genitiv)}
    \item Sie hilft \underline{\hspace{1.5cm}} alten Mann. \quad \textit{(Lösung: dem -- Dativ)}
    \item \underline{\hspace{1.5cm}} Frau gab mir \underline{\hspace{1.5cm}} Buch. \quad \textit{(Lösung: Die, ein -- Nominativ, Akkusativ)}
    \item Wir denken an \underline{\hspace{1.5cm}} Zukunft. \quad \textit{(Lösung: die -- Akkusativ)}
    \item \underline{\hspace{1.5cm}} Kindern geht es gut. \quad \textit{(Lösung: Den -- Dativ)}
    \item Trotz \underline{\hspace{1.5cm}} Regens gingen wir los. \quad \textit{(Lösung: des -- Genitiv)}
    \item Ich habe \underline{\hspace{1.5cm}} Apfel gegessen. \quad \textit{(Lösung: einen -- Akkusativ)}
    \item Er spricht mit \underline{\hspace{1.5cm}} Lehrerin. \quad \textit{(Lösung: der -- Dativ)}
    \item Das ist das Haus \underline{\hspace{1.5cm}} Freundes. \quad \textit{(Lösung: des -- Genitiv)}
    \item Für \underline{\hspace{1.5cm}} gute Arbeit danke ich dir. \quad \textit{(Lösung: die -- Akkusativ)}
    \item \underline{\hspace{1.5cm}} Studenten lernen fleißig. \quad \textit{(Lösung: Die -- Nominativ)}
    \item Während \underline{\hspace{1.5cm}} Pause tranken wir Kaffee. \quad \textit{(Lösung: der -- Genitiv)}
    \item Ich gebe \underline{\hspace{1.5cm}} Mädchen ein Geschenk. \quad \textit{(Lösung: dem -- Dativ)}
    \item \underline{\hspace{1.5cm}} Berg ist hoch. \quad \textit{(Lösung: Der -- Nominativ)}
    \item Wir fahren nach \underline{\hspace{1.5cm}} Schweiz. \quad \textit{(Lösung: der -- Dativ/Akkusativ fix)}
    \item Ohne \underline{\hspace{1.5cm}} Geld kannst du nichts kaufen. \quad \textit{(Lösung: - -- Nullartikel)}
    \item \underline{\hspace{1.5cm}} Eltern sind stolz. \quad \textit{(Lösung: Die -- Nominativ)}
    \item Ich erinnere mich an \underline{\hspace{1.5cm}} Sommer. \quad \textit{(Lösung: den -- Akkusativ)}
    \item \underline{\hspace{1.5cm}} Probleme sind gelöst. \quad \textit{(Lösung: Die -- Nominativ)}
    \item Das ist der Wunsch \underline{\hspace{1.5cm}} Königs. \quad \textit{(Lösung: des -- Genitiv)}
    \item Der Rat \underline{\hspace{1.5cm}} Freundes war wichtig. \quad \textit{(Lösung: des -- Genitiv)}
    \item Ich gedenke \underline{\hspace{1.5cm}} Helden. \quad \textit{(Lösung: der -- Genitiv)}
    \item Er bedarf \underline{\hspace{1.5cm}} Hilfe. \quad \textit{(Lösung: der -- Genitiv)}
    \item Sie wurde \underline{\hspace{1.5cm}} Lage Herr. \quad \textit{(Lösung: der -- Genitiv)}
    \item Die Lösung \underline{\hspace{1.5cm}} Problems war einfach. \quad \textit{(Lösung: des -- Genitiv)}
    \item Der Erfolg \underline{\hspace{1.5cm}} Projekts freut uns. \quad \textit{(Lösung: des -- Genitiv)}
    \item Das liegt am Wetter \underline{\hspace{1.5cm}} Woche. \quad \textit{(Lösung: der -- Genitiv)}
    \item Der Chef \underline{\hspace{1.5cm}} Abteilung ist streng. \quad \textit{(Lösung: der -- Genitiv)}
    \item Die Höhe \underline{\hspace{1.5cm}} Berge ist impressive. \quad \textit{(Lösung: der -- Genitiv)}
\end{enumerate}

\section{Erweiterte Übungen: Kasusbestimmung}
\label{sec:ErweiterteUebungen}

Bestimmen Sie den Kasus der unterstrichenen Nomen:

\begin{enumerate}
    \item \underline{\hspace{1.5cm}} Junge hat Hunger. (Nominativ)
    \item Ich sehe \underline{\hspace{1.5cm}} junge Mann. (Akkusativ)
    \item Sie gibt \underline{\hspace{1.5cm}} Kind ein Buch. (Dativ)
    \item Das Auto \underline{\hspace{1.5cm}} Nachbarin ist rot. (Genitiv)
    \item Wir helfen \underline{\hspace{1.5cm}} alten Frau. (Dativ)
    \item Der Lehrer erklärt \underline{\hspace{1.5cm}} Schüler die Regel. (Dativ, Akkusativ)
    \item Ich erinnere mich an \underline{\hspace{1.5cm}} schöne Tag. (Akkusativ)
    \item Trotz \underline{\hspace{1.5cm}} Krankheit ging er zur Arbeit. (Genitiv)
    \item Der Preis \underline{\hspace{1.5cm}} Reise war hoch. (Genitiv)
    \item Sie ist stolz auf \underline{\hspace{1.5cm}} Erfolg. (Akkusativ)
    \item Das ist das Buch \underline{\hspace{1.5cm}} Autors. (Genitiv)
    \item Ich vertraue \underline{\hspace{1.5cm}} Freund. (Dativ)
    \item Der Brief \underline{\hspace{1.5cm}} Schwester liegt auf dem Tisch. (Genitiv)
    \item Sie freut sich über \underline{\hspace{1.5cm}} Geschenk. (Akkusativ)
    \item Wegen \underline{\hspace{1.5cm}} Staubs fiel der Unterricht aus. (Genitiv)
\end{enumerate}

\section{Regionale Besonderheiten}
\label{sec:RegionaleBesonderheiten}

\subsection{Genitiv-Ersatz durch von + Dativ}
In der gesprochenen Umgangssprache wird der Genitiv oft durch „von + Dativ" ersetzt:
\begin{itemize}
    \item \textit{Standard}: Das Auto des Mannes
    \item \textit{Umgangssprache}: Das Auto von dem Mann / vom Mann
    \item \textit{Standard}: Das Buch des Kindes
    \item \textit{Umgangssprache}: Das Buch von dem Kind
\end{itemize}

Dies gilt jedoch als grammatikalisch inkorrekt in der Standardsprache und im schriftlichen Gebrauch.

\section{Lösungen zu den Übungen}
\label{sec:Loesungen}

Die Lösungen zu allen Übungen finden Sie im Anhang.

\chapterend